%=========================================================
% C++ Chapter 13 Lab Handout (Verbatim Korean-safe)
% Topic: Exception Handling & C Linkage (extern "C")
% Compile with XeLaTeX
%=========================================================
\documentclass[11pt,a4paper]{article}

\usepackage[margin=1in]{geometry}
\usepackage{fontspec}
\usepackage{kotex}
\usepackage{hyperref}
\usepackage{enumitem}
\usepackage{fancyhdr}
\usepackage{titlesec}
\usepackage{xcolor}
\usepackage{tcolorbox}
\usepackage{fvextra}

% ---------- Fonts ----------
\setmonofont{D2Coding}[Scale=MatchLowercase]

% ---------- Colors ----------
\definecolor{CodeBg}{RGB}{248,248,248}
\definecolor{CodeFrame}{RGB}{220,220,220}
\definecolor{Accent}{RGB}{40,80,160}

% ---------- Header ----------
\pagestyle{fancy}
\fancyhf{}
\lhead{\textbf{C++ 13장 실습}}
\rhead{\textbf{예외 처리와 C 링크}}
\cfoot{\thepage}

\titleformat{\section}{\Large\bfseries\color{Accent}}{\thesection}{0.6em}{}
\titleformat{\subsection}{\large\bfseries}{\thesubsection}{0.6em}{}

% ---------- Boxes ----------
\newtcolorbox{GoalBox}{
  colback=CodeBg,
  colframe=CodeFrame,
  arc=3pt,
  boxrule=0.8pt,
  left=10pt,right=10pt,top=8pt,bottom=8pt
}
\newcommand{\filetag}[1]{\texttt{\color{Accent}#1}}

% ---------- Code block environment (Unicode-safe) ----------
% IMPORTANT: no commandchars => C++ braces {} print correctly.
\DefineVerbatimEnvironment{cppcode}{Verbatim}{
  fontsize=\small,
  frame=single,
  framerule=0.6pt,
  rulecolor=\color{CodeFrame},
  numbers=left,
  numbersep=8pt,
  baselinestretch=1,
  breaklines=true,
  breaksymbolleft=,
  breaksymbolright=,
  xleftmargin=1.5em
}

\begin{document}

\begin{center}
  {\LARGE \textbf{C++ 13장 실습: 예외 처리와 C 언어와의 링크 지정}}\\
  \vspace{6pt}
  {\large (try-throw-catch, exception class, nested try, C++ name mangling, extern "C")}
\end{center}

\vspace{8pt}

\section*{학습 목표}
\begin{GoalBox}
\begin{enumerate}[leftmargin=18pt]
  \item 실행 오류와 오류 처리의 일반적인 방법을 복습한다.
  \item 예외와 예외 처리의 개념을 이해한다.
  \item try-throw-catch로 구성되는 예외 처리의 기본 형식을 안다.
  \item 예외 클래스를 작성하여 예외를 처리하는 방법을 안다.
  \item C++의 이름 규칙을 이해한다.
  \item C++ 코드와 C 코드의 링킹에 이름 규칙의 중요성을 이해한다.
  \item extern "C"를 이용하여 C++ 코드와 C 코드의 성공적인 링킹 방법을 안다.
\end{enumerate}
\end{GoalBox}

\section*{실습 운영 규칙}
\begin{itemize}[leftmargin=18pt]
  \item 모든 예제는 \texttt{main()} 포함, \textbf{독립 실행}입니다.
  \item 예외 처리 실습의 핵심은 \textbf{정상 흐름(정상 값 계산)}과 \textbf{비정상 흐름(예외 경로)}를 분리하는 것입니다.
  \item \textbf{주의:} 슬라이드에 등장하는 \texttt{throw(int)} 같은 ``예외 명세(exception specification)''는 현대 C++에서 사용을 권장하지 않습니다(구형 문법). 본 실습에서는 \textbf{개념만 이해}하고, 코드는 쓰지 않는 방식으로 제공합니다.
  \item extern "C"는 ``C로 컴파일하라''가 아니라, \textbf{이 함수의 링크 이름을 C 규칙으로 만들라}는 의미입니다.
\end{itemize}

\tableofcontents
\newpage

%=========================================================
\section{실행 오류 복습: 예외 상황 대응이 없는 코드 (예제 13-1)}
\begin{GoalBox}
\textbf{핵심}
\begin{itemize}[leftmargin=18pt]
  \item 음수 지수 같은 ``예상 밖 입력''이 들어오면 결과가 엉뚱해질 수 있음.
  \item 오류를 발견했는데도 계속 계산하면, \textbf{틀린 값이 정상처럼 출력}되는 것이 가장 위험.
\end{itemize}
\end{GoalBox}

\noindent\filetag{L1\_getExp\_no\_handling.cpp}
\begin{cppcode}
#include <iostream>
using namespace std;

// base의 exp 지수승을 계산하여 리턴 (하지만 exp<0을 고려하지 않음!)
int getExp(int base, int exp) {
    int value = 1;
    for(int n=0; n<exp; n++) // exp가 음수면 루프가 0번 실행됨 -> value==1
        value *= base;
    return value;
}

int main() {
    cout << "2의 3승 = " << getExp(2, 3) << "\n";
    cout << "2의 -3승 = " << getExp(2, -3) << "\n"; // 잘못된 결과가 조용히 출력됨
}
\end{cppcode}

%=========================================================
\section{전통적 오류 처리 1: 리턴 값으로 오류 표시 (예제 13-2)}
\begin{GoalBox}
\textbf{핵심}
\begin{itemize}[leftmargin=18pt]
  \item 오류를 \texttt{-1} 같은 값으로 표시하면, \textbf{정상 결과와 충돌}할 수 있음.
  \item 호출자가 매번 \texttt{if (ret == -1)} 같은 코드를 반복해야 하므로 ``번거롭고 실수''가 생김.
\end{itemize}
\end{GoalBox}

\noindent\filetag{L2\_getExp\_return\_error.cpp}
\begin{cppcode}
#include <iostream>
using namespace std;

int getExp(int base, int exp) {
    if(base <= 0 || exp <= 0) return -1; // 오류를 -1로 표시(취약)
    int value = 1;
    for(int n=0; n<exp; n++) value *= base;
    return value;
}

int main() {
    int v = getExp(2, 3);
    if(v != -1) cout << v << "\n";
    else        cout << "오류\n";

    int e = getExp(2, -3);
    if(e != -1) cout << e << "\n";
    else        cout << "오류\n";
}
\end{cppcode}

%=========================================================
\section{전통적 오류 처리 2: bool + 참조 매개변수(out) (예제 13-3)}
\begin{GoalBox}
\textbf{핵심}
\begin{itemize}[leftmargin=18pt]
  \item 리턴 값은 \textbf{성공/실패}만 전달하고, 계산 결과는 \textbf{참조 매개변수}로 전달.
  \item 리턴 값 의미가 단일화되어(오류 상태만 표시) 비교적 깔끔해짐.
\end{itemize}
\end{GoalBox}

\noindent\filetag{L3\_getExp\_bool\_out.cpp}
\begin{cppcode}
#include <iostream>
using namespace std;

bool getExp(int base, int exp, int& out) {
    if(base <= 0 || exp <= 0) return false;
    int value = 1;
    for(int n=0; n<exp; n++) value *= base;
    out = value;
    return true;
}

int main() {
    int v;
    if(getExp(2, 3, v)) cout << v << "\n";
    else cout << "오류\n";

    int e;
    if(getExp(2, -3, e)) cout << e << "\n";
    else cout << "오류\n";
}
\end{cppcode}

%=========================================================
\section{예외 처리 기본 형식: try-throw-catch (슬라이드 핵심)}
\begin{GoalBox}
\textbf{핵심}
\begin{itemize}[leftmargin=18pt]
  \item try: 예외가 발생할 가능성이 있는 코드를 묶는다.
  \item throw: 예외 발생을 알리고, \textbf{정상 흐름을 즉시 중단}하고 catch로 점프.
  \item catch: 예외 타입에 맞는 블록에서 처리한다(여러 catch 가능).
\end{itemize}
\end{GoalBox}

%=========================================================
\section{예제 13-4: 0으로 나누기 예외 처리(입력 검증 + throw int)}
\begin{GoalBox}
\textbf{핵심}
\begin{itemize}[leftmargin=18pt]
  \item 입력값 n이 0/음수면 계산을 하지 않고 throw로 예외 경로로 보낸다.
  \item catch에서 메시지를 출력하고, 루프를 계속 돌릴지(continue), 종료할지 결정.
\end{itemize}
\end{GoalBox}

\noindent\filetag{L4\_average\_divide\_exception.cpp}
\begin{cppcode}
#include <iostream>
using namespace std;

int main() {
    int n, sum, average;
    while(true) {
        cout << "합을 입력하세요>>";
        cin >> sum;
        cout << "인원수를 입력하세요>>";
        cin >> n;

        try {
            if(n <= 0) throw n;          // 예외 발생
            average = sum / n;           // 정상 계산
        }
        catch(int x) {
            cout << "예외 발생!! " << x << "으로 나눌 수 없음\n\n";
            continue;                    // 다음 입력으로
        }

        cout << "평균 = " << average << "\n\n";
    }
}
\end{cppcode}

%=========================================================
\section{함수에서 예외 던지기: getExp 완결판(예제 13-5)}
\begin{GoalBox}
\textbf{핵심}
\begin{itemize}[leftmargin=18pt]
  \item ``오류가 나면 리턴''이 아니라, \textbf{오류를 발견한 곳에서 throw}한다.
  \item 호출자는 try/catch에서 \textbf{한 곳에 모아 처리}할 수 있음.
\end{itemize}
\end{GoalBox}

\noindent\filetag{L5\_getExp\_throw.cpp}
\begin{cppcode}
#include <iostream>
using namespace std;

int getExp(int base, int exp) {
    if(base <= 0 || exp <= 0) throw "음수 사용 불가";
    int value = 1;
    for(int n=0; n<exp; n++) value *= base;
    return value;
}

int main() {
    try {
        cout << "2^3 = " << getExp(2, 3) << "\n";
        cout << "2^-3 = " << getExp(2, -3) << "\n"; // 여기서 throw
        cout << "이 줄은 실행되지 않음\n";
    }
    catch(const char* msg) {
        cout << "예외 발생!! " << msg << "\n";
    }
}
\end{cppcode}

%=========================================================
\section{예제 13-6: 문자열을 정수로 변환 (문자 발견 시 throw)}
\begin{GoalBox}
\textbf{핵심}
\begin{itemize}[leftmargin=18pt]
  \item 변환 실패 시, 실패 원인(입력 문자열)을 예외로 던질 수 있다.
  \item catch에서 사용자에게 어떤 입력이 문제였는지 알려줄 수 있다.
\end{itemize}
\end{GoalBox}

\noindent\filetag{L6\_stringToInt.cpp}
\begin{cppcode}
#include <iostream>
#include <cstring>
using namespace std;

int stringToInt(const char x[]) {
    int sum = 0;
    int len = strlen(x);
    for(int i=0; i<len; i++) {
        if(x[i] >= '0' && x[i] <= '9')
            sum = sum*10 + (x[i]-'0');
        else
            throw x; // 변환 불가 -> 입력 문자열을 예외로 던짐
    }
    return sum;
}

int main() {
    try {
        cout << "\"123\" -> " << stringToInt("123") << "\n";
        cout << "\"1A3\" -> " << stringToInt("1A3") << "\n";
    }
    catch(const char* s) {
        cout << s << " 처리에서 예외 발생!!\n";
    }
}
\end{cppcode}

%=========================================================
\section{예제 13-7: 예외 처리를 가진 스택 클래스}
\begin{GoalBox}
\textbf{핵심}
\begin{itemize}[leftmargin=18pt]
  \item 스택이 full/empty면 문자열 예외를 던져 호출자에게 알린다.
  \item 호출자는 try/catch로 한 번에 처리할 수 있다.
\end{itemize}
\end{GoalBox}

\noindent\filetag{L7\_MyStack\_exceptions.cpp}
\begin{cppcode}
#include <iostream>
using namespace std;

class MyStack {
    int data[3]; // 실습용으로 작게(overflow를 쉽게 재현)
    int tos;
public:
    MyStack(): tos(-1) {}
    void push(int n) {
        if(tos == 2) throw "Stack Full";
        data[++tos] = n;
    }
    int pop() {
        if(tos == -1) throw "Stack Empty";
        return data[tos--];
    }
};

int main() {
    MyStack st;
    try {
        st.push(100);
        st.push(200);
        st.push(300);
        // st.push(400); // (실습) 주석 풀면 "Stack Full"
        cout << st.pop() << "\n";
        cout << st.pop() << "\n";
        cout << st.pop() << "\n";
        cout << st.pop() << "\n"; // Stack Empty
    }
    catch(const char* msg) {
        cout << "예외 발생: " << msg << "\n";
    }
}
\end{cppcode}

%=========================================================
\section{중첩 try: 안쪽 예외는 안쪽 catch가 우선(슬라이드)}
\begin{GoalBox}
\textbf{핵심}
\begin{itemize}[leftmargin=18pt]
  \item try 안에 try를 중첩할 수 있다.
  \item 예외는 \textbf{가장 가까운(안쪽) catch}부터 매칭을 시도한다.
\end{itemize}
\end{GoalBox}

\noindent\filetag{L8\_nested\_try.cpp}
\begin{cppcode}
#include <iostream>
using namespace std;

int main() {
    try {
        throw 3; // 바깥 예외

        try {
            throw "abc"; // 안쪽 예외
            throw 5;
        }
        catch(int inner) { // 안쪽 try에 연결된 catch
            cout << "inner int: " << inner << "\n";
        }
    }
    catch(const char* s) { // 바깥 try의 catch
        cout << "outer const char*: " << s << "\n";
    }
    catch(int outer) {
        cout << "outer int: " << outer << "\n";
    }
}
\end{cppcode}

%=========================================================
\section{예외 클래스 만들기(예제 13-8 핵심 재구성)}
\begin{GoalBox}
\textbf{핵심}
\begin{itemize}[leftmargin=18pt]
  \item 예외를 객체로 던지면, 라인/함수/메시지 등 \textbf{풍부한 정보}를 담을 수 있다.
  \item 공통 기반 예외(MyException)를 만들고, 구체 예외를 상속으로 분리 가능.
\end{itemize}
\end{GoalBox}

\noindent\filetag{L9\_exception\_classes.cpp}
\begin{cppcode}
#include <iostream>
#include <string>
using namespace std;

class MyException {
    int lineNo;
    string func;
    string msg;
public:
    MyException(int n, string f, string m) : lineNo(n), func(f), msg(m) {}
    void print() const { cout << func << ":" << lineNo << ", " << msg << "\n"; }
};

class DivideByZeroException : public MyException {
public:
    DivideByZeroException(int n, string f, string m) : MyException(n,f,m) {}
};

class InvalidInputException : public MyException {
public:
    InvalidInputException(int n, string f, string m) : MyException(n,f,m) {}
};

int main() {
    int x, y;
    try {
        cout << "나눗셈: 양의 정수 2개 입력>> ";
        cin >> x >> y;

        if(x < 0 || y < 0)
            throw InvalidInputException(100, "main()", "음수 입력 예외 발생");
        if(y == 0)
            throw DivideByZeroException(103, "main()", "0으로 나누는 예외 발생");

        cout << (double)x / (double)y << "\n";
    }
    catch(const DivideByZeroException& e) { e.print(); }
    catch(const InvalidInputException& e)  { e.print(); }
}
\end{cppcode}

%=========================================================
\section{C/C++ 링크: 이름 맹글링(name mangling)과 extern "C"}
\begin{GoalBox}
\textbf{핵심}
\begin{itemize}[leftmargin=18pt]
  \item C는 함수 이름이 단순(\_f)하지만, C++은 오버로딩 때문에 매개변수 타입까지 포함한 ``복잡한 이름''을 만들 수 있다.
  \item 그래서 C로 컴파일된 f.c의 심볼(\_f)을, C++에서 기대하는 심볼(?f@@...)과 매칭 못해 링크 에러가 발생할 수 있다.
  \item 해결: C로 컴파일된 함수를 C++에서 선언할 때 \texttt{extern "C"}로 감싸 \textbf{C 링크 이름 규칙}을 사용하게 한다.
\end{itemize}
\end{GoalBox}

\subsection*{extern "C" 사용 예(단일 함수)}
\noindent\filetag{L10\_externC\_single.cpp}
\begin{cppcode}
/*
(개념 실습)
- f.c 는 C 컴파일러로 컴파일된다고 가정한다.
- main.cpp 에서 f()를 extern "C"로 선언하면 링크 성공.
*/

/* main.cpp */
#include <iostream>
using namespace std;

extern "C" int f(int x, int y); // C 링킹 규칙 사용
int g(int x, int y);            // C++ 링킹 규칙 사용

int main() {
    cout << f(2,5) << "\n";
    cout << g(2,5) << "\n";
}

/* g.cpp (C++) */
int g(int x, int y) { return x - y; }

/* f.c (C) */
int f(int x, int y) { return x + y; }
\end{cppcode}

\subsection*{extern "C" 사용 예(여러 함수/헤더)}
\noindent\filetag{L10b\_externC\_header.cpp}
\begin{cppcode}
/*
extern "C" {
    int f(int x, int y);
    void g();
    char s(int []);
}

extern "C" {
#include "mycfunction.h"
}
*/
\end{cppcode}

%=========================================================
\section*{마무리 체크리스트}
\begin{GoalBox}
\begin{itemize}[leftmargin=18pt]
  \item 리턴 값 오류 처리의 한계를 설명할 수 있는가?
  \item bool+out 패턴의 장단점을 말할 수 있는가?
  \item try/throw/catch의 제어 흐름(throw 이후 코드는 실행 안 됨)을 설명할 수 있는가?
  \item 함수에서 예외를 던지고, 호출자에서 한 번에 처리하는 구조를 작성할 수 있는가?
  \item 예외 클래스로 라인/함수/메시지를 담아 처리할 수 있는가?
  \item 이름 맹글링이 왜 필요한지(오버로딩) 설명할 수 있는가?
  \item C 함수 링킹 실패 원인과 extern "C" 해결책을 말할 수 있는가?
\end{itemize}
\end{GoalBox}

\end{document}
